\begin{center}\begin{tikzpicture}[font=\small,>=stealth,line join=round]
\draw[gray!40,rounded corners=3pt] (0,0) rectangle (14.2,8.3);\node at (7.1,7.700000000000001) {例7：从原函数里分出偶部分与奇部分};
\begin{scope}[shift={(4.1,3.102631578947368)},x=2.0cm,y=0.2421052631578947cm]\draw[->,gray] (-1.5,0)--(1.5,0) node[right]{$x$};\draw[->,gray] (0,-6)--(0,13) node[above]{$y$};\draw[blue,thick,solid,domain=-1.2:1.2,samples=180] plot(\x,{(\x)^3+2*(\x)^2+3*\x+4});\draw[violet,thick,dashed,domain=-1.2:1.2,samples=180] plot(\x,{2*(\x)^2+4});\draw[orange!85!black,thick,dotted,domain=-1.2:1.2,samples=180] plot(\x,{(\x)^3+3*\x});\end{scope}\node[anchor=west,align=left,text width=5.0cm] at (8.4,4.0) {蓝：原函数 $H$\\[10pt]紫：$E=2x^2+4$\\[10pt]橙：$O=x^3+3x$\\[10pt]任意 $x$ 处都有 $H=E+O$};
\node[text width=12.6cm,align=center] at (7.1,.65) {分解之后，左右相同的高度归入紫线，左右相反的高度归入橙线。};
\end{tikzpicture}\end{center}
